\section*{Roadmap and Reading Guidance}
\addcontentsline{toc}{section}{Roadmap and Reading Guidance}

The title names three pillars, each developed in its own parts. The architecture---one
closed system, the atomic move, conservation, and the three-home state model---is
Parts~I--II (\S\S1--4). The instrument lifecycle---valuation, contracts, deterministic and
optional lifecycle events, and the generalisation to obligations and securities
lending---is Parts~III and~V (\S\S5--8, \S14, \S15). CDM integration, with accounts,
reporting, and the settlement interface, is Part~IV (\S\S9--13). Part~VI (\S16, \S\S17--20)
states the formal properties, limitations, and context.

The reading paths below give, for each audience, the shortest route to a working
understanding, what to add when more depth is required, and what carries no cost on a first
pass.

\begin{longtable}{@{}p{0.16\textwidth}p{0.30\textwidth}p{0.28\textwidth}p{0.18\textwidth}@{}}
\toprule
\textbf{Audience} & \textbf{Minimal path} & \textbf{Then if needed} & \textbf{May skip on first pass} \\
\midrule
\endfirsthead
\toprule
\textbf{Audience} & \textbf{Minimal path} & \textbf{Then if needed} & \textbf{May skip on first pass} \\
\midrule
\endhead
\bottomrule
\endfoot
Regulator &
Preface, \S1, the conservation core of \S2, \S17, \S18 &
\S12, \S19, \S20, App.~C, App.~H, App.~I &
the Haskell, \S4 construction detail, \S6, \S8 detail, \S\S14--16 internals \\
\addlinespace
Implementer &
\S2, \S3, \S4, \S6, \S7, \S11, \S16 &
\S8, \S12, \S13 + App.~F, \S14, \S15, App.~A/B/G &
\S5 profit-and-loss derivations, \S9, \S10, \S17, worked examples \\
\addlinespace
Quant &
\S2, \S5, \S10, App.~D &
\S8, \S15 &
\S3, \S11, \S12, \S\S13--14 mechanics, \S17 \\
\addlinespace
Risk manager &
\S2, \S5, \S7, \S9, \S10, \S15 &
\S17, App.~C &
\S3, \S6 internals, \S11, \S\S13--14, the Haskell \\
\bottomrule
\end{longtable}

The progression from scalar to vector is deliberate. The scalar wallet balance carries the
whole argument through Parts~I--IV; the six-coordinate position vector and securities
lending wait until Part~V, introduced only once obligations are in hand.
